\subsection{自反模的局部构造}

沿用第 2 节的记号和假设。若某性质满足``对几乎所有 \(p \in P\)''，即不满足该性质的 \(p \in P\) 所成集合是有限集，则称该性质成立。

定理 3. 设 V 是 K 上的有限秩向量空间，M 是 V 关于 A 的格。

\begin{enumerate}
\def\labelenumi{(\roman{enumi})}
\item
  设 N 是 V 关于 A 的格；则对 A 的每个素理想 p，\(N_{p}\) 是 V 关于 \(A_{p}\) 的格，并且对几乎所有 \(p \in P\)，有 \(N_{p} = M_{p}\)。
\item
  反之，设对每个 \(\mathfrak{p} \in \mathbf{P}\) 给定一个 \(\mathrm{N}(\mathfrak{p})\)\(\mathbf{V}\) 关于 \(\mathbf{A}_{\mathfrak{p}}\) 的格 ，使得对几乎所有  有 \(\mathrm{N}(\mathfrak{p}) = \mathrm{M}_{\mathfrak{p}}\)。则 \(\mathfrak{p} \in \mathbf{P}\)\(\mathrm{N} = \bigcap_{\mathfrak{p} \in \mathbf{P}} \mathrm{N}(\mathfrak{p})\) 是 \(\mathbf{A}\) 关于 \(\mathbf{A}\) 的自反格，并且它是唯一满足对所有 \(\mathrm{N}'\) 有  的 \(\mathbf{V}\) 关于 \(\mathbf{A}\) 的自反格 \(\mathrm{N}_{\mathfrak{p}}' = \mathrm{N}(\mathfrak{p})\)\(\mathfrak{p} \in \mathbf{P}\)。
\item
  第一断言由第 1 节命题 4 得出。此外，存在  中的 \(x, y\)，使得 \(\mathbf{K}^*\)\(x\mathrm{N} \subset \mathrm{M} \subset y\mathrm{N}\)（第 1 节，命题 2）；我们知道，对几乎所有 \(\mathfrak{p} \in \mathbf{P}\)，有 \(v_{\mathfrak{p}}(x) = v_{\mathfrak{p}}(y) = 0\)（§ 1，第 6 节，定理 4），这表明 x 和 y 在 \(x\)\(y\)\(\mathbf{A}_{\mathfrak{p}}\) 中可逆，从而 \(\mathbf{M}_{\mathfrak{p}} = \mathbf{N}_{\mathfrak{p}}\)。
\item
  可以将 M 换为 \(x^{-1}\mathbf{M}\)，其中 \(x \neq 0\) 属于 A，并假设对所有  都有 \(\mathrm{N}(\mathfrak{p}) \subset \mathbf{M}_{\mathfrak{p}}\)。令 \(\mathfrak{p} \in \mathrm{P}\)\(\mathfrak{p}_1, \ldots, \mathfrak{p}_h\) 为 P 中满足当 \(\mathrm{N}(\mathfrak{p}) = \mathrm{M}_{\mathfrak{p}}\)\(\mathfrak{p}\) 不同于各个 \(\mathfrak{p}_i\) 时有  的元素（\(1 \leqslant i \leqslant h\)）；记作：
\end{enumerate}

\[
Q = M \cap N (p _ {1}) \cap \dots \cap N (p _ {h}).
\]

由于每个 \(\mathrm{N}(p_{i})\) 都包含一个关于 \(A_{p_{i}}\) 的自由格，所以它当然包含一个关于 A 的 V 的格；因而 Q 包含一个关于 A 的 V 的格（第 1 节，命题 3），又由于 Q 包含于 M，Q 是关于 A 的格。为证明对所有  有 \(Q_{p} = N(p)\)，我们将使用下列引理：\(p \in P\)

引理 1. 设 \(\mathfrak{p}\) 和 \(\mathfrak{p}'\) 是 \(\mathbf{A}\) 的两个素理想，且 (0) 是包含于 \(\mathbf{A}\)\(\mathfrak{p} \cap \mathfrak{p}'\) 的 \(\mathbf{A}\) 的唯一素理想。则对于  的每个子--模 \(\mathbf{E}\)，有 \(\mathbf{V}\)\((\mathbf{E}_{\mathfrak{p}})_{\mathfrak{p}'} = \mathbf{K} \cdot \mathbf{E}\)。

令 S 为 A 的乘性子集 \((\mathbf{A}-\mathfrak{p})(\mathbf{A}-\mathfrak{p}')\)；由第二章，§2，第 3 节，命题 7，\((\mathrm{E}_{\mathfrak{p}})_{\mathfrak{p}'}=\mathrm{S}^{-1}\mathrm{E}\)。此外，\(A\subset S^{-1}A\subset K\)；\(S^{-1}A\) 的素理想对应于 A 中满足 \(q\cap S=\varnothing\) 的素理想 q（第二章，§2，第 5 节，命题 11），而根据假设，(0) 是 A 中唯一不与 S 相交的素理想；因此 \(S^{-1}A=K\)，且 \(S^{-1}E=K.E\)。

现在回到 (ii) 的证明。若 \(p \in P\) 不同于 \(p_i (1 \leqslant i \leqslant h)\)，将引理 1 应用于 \(\mathrm{N}(\mathfrak{p}_i)\)，得到 \((\mathrm{N}(\mathfrak{p}_i))_{\mathfrak{p}} = ((\mathrm{N}(\mathfrak{p}_i))_{\mathfrak{p}_i})_{\mathfrak{p}} = \mathrm{K}. \mathrm{N}(\mathfrak{p}_i) = \mathrm{V}\)，因为 \(p_i\) 和 p 的高度均为 1。于是

\[
\mathbf {Q} _ {\mathfrak {p}} = \mathbf {M} _ {\mathfrak {p}} \cap (\mathbf {N} (\mathfrak {p} _ {1})) _ {\mathfrak {p}} \cap \dots \cap (\mathbf {N} (\mathfrak {p} _ {h})) _ {\mathfrak {p}} = \mathbf {M} _ {\mathfrak {p}} = \mathbf {N} (\mathfrak {p})
\]

（第二章，§ 2，第 4 节）。另一方面，若 \(\mathfrak{p}\) 等于 \(\mathfrak{p}_i\)（\(1 \leqslant i \leqslant h\)），则由上述论证，当  时 \((\mathrm{N}(\mathfrak{p}_i))_{\mathfrak{p}_j} = \mathrm{V}\)，且 \(i \neq j\)\((\mathrm{N}(\mathfrak{p}_i))_{\mathfrak{p}_i} = \mathrm{N}(\mathfrak{p}_i)\)，从而

\[
\mathbf {Q} _ {\mathfrak {p} _ {i}} = \mathbf {M} _ {\mathfrak {p} _ {i}} \cap \mathbf {N} (\mathfrak {p} _ {i}) = \mathbf {N} (\mathfrak {p} _ {i}).
\]

因此我们已经证明，对所有 ，有 \(Q_{\mathfrak{p}} = N(\mathfrak{p})\)。于是\(\mathfrak{p} \in P\)

\[
\mathbf {N} = \mathbf {Q} ^ {* *} = \bigcap_ {\mathfrak {p} \in \mathbf {P}} \mathbf {Q} _ {\mathfrak {p}}
\]

是自反的，并满足对所有  有 \(N_{p} = Q_{p} = N(p)\)；唯一性立即由第 2 节的定理 2 得出。\(p \in P\)

注。设 L 是 V 关于 A 的自由格。由于对 ，\(A_{p}\) 是主理想整环，\(p \in P\)\(N(p)\) 是秩与 L 相同的自由 \(A_{p}\)-模，并且存在 \(u(p) \in \mathbf{GL}(V)\)，使得 \(u(p)(L_{p}) = N_{p}\)；此外，这个条件把 \(u(p)\) 确定到右乘 \(\mathbf{GL}(L_{p})\) 中的一个元素为止。条件“对几乎所有  有 \(N(p) = L_{p}\)”意味着必有 \(p \in P\)\(u(p) \in \mathbf{GL}(L_{p})\)

对几乎所有 \(\mathfrak{p} \in \mathbb{P}\) 都如此。满足后一性质的族 \((u(\mathfrak{p}))_{\mathfrak{p} \in \mathbb{P}}\) 构成一个乘法群 \(\mathbf{GL}_a(V)\)，它包含乘积 \(\prod_{\mathfrak{p} \in \mathbb{P}} \mathbf{GL}(L_{\mathfrak{p}})\) 作为子群。于是定理 3 表明，V 的自反格集合与齐性空间 \(V\)\(\mathbf{GL}_a(V)/\prod_{\mathfrak{p} \in \mathbb{P}} \mathbf{GL}(L_{\mathfrak{p}})\) 存在典范一一对应。若选取 L 在 A 上的基 \((e_i)_{1 \leq i \leq n}\)，则 \(L\)\(A\)\(\mathbf{GL}(V)\)（相应地 \(\mathbf{GL}(L_{\mathfrak{p}})\)）可识别为可逆矩阵群 \(\mathbf{GL}(n, K)\)（相应地 \(\mathbf{GL}(n, A_{\mathfrak{p}})\)），而群 \(\mathbf{GL}_a(V)\) 可识别为阶为 \(n\) 的矩阵系群 \((U(\mathfrak{p}))_{\mathfrak{p} \in P}\)，满足对所有  有 \(U(\mathfrak{p}) \in \mathbf{GL}(n, K)\)，并且对几乎所有 \(\mathfrak{p} \in P\) 有 \(U(\mathfrak{p}) \in \mathbf{GL}(n, A_{\mathfrak{p}})\)。若 A 是 Dedekind 整环，则群 \(\mathfrak{p} \in P\)\(A\)\(\mathbf{GL}_a(V)\) 也可识别为群 \(\mathbf{GL}(n, A)\)，其中 A 是限制阿代尔环（§ 2，第 4 节）。\(A\)

